% =============================================================================
%  Differentiable Inverse Reachability Operator for RM75
%  Complete mathematical exposition of the IRDPLAYGROUND design
% =============================================================================
% Overleaf / local: Compiler = pdfLaTeX
%   pdflatex ird_operator_mathematics.tex
% Do NOT use fontspec / Noto CJK here (English-only).
% Chinese version: ird_operator_mathematics_zh.tex  (XeLaTeX + ctex)
% =============================================================================
\documentclass[11pt,a4paper]{article}

\usepackage[margin=2.4cm]{geometry}
\usepackage{amsmath,amssymb,amsthm,mathtools}
\usepackage{bm}
\usepackage{booktabs}
\usepackage{enumitem}
\usepackage{hyperref}
\usepackage{microtype}
\usepackage[T1]{fontenc}
\usepackage{lmodern}
\usepackage{xcolor}
\usepackage{listings}
\usepackage{graphicx}
\usepackage{cleveref}

\hypersetup{
  colorlinks=true,
  linkcolor=blue!50!black,
  citecolor=blue!50!black,
  urlcolor=blue!60!black
}

\newtheorem{definition}{Definition}
\newtheorem{remark}{Remark}
\newtheorem{proposition}{Proposition}

\newcommand{\SE}{\mathrm{SE}}
\newcommand{\SO}{\mathrm{SO}}
\newcommand{\R}{\mathbb{R}}
\newcommand{\Tcp}{T_{\mathrm{tcp}}}
\newcommand{\Taxis}{T_{\mathrm{axis}}}
\newcommand{\Tfl}{T_{\mathrm{flange}}}
\newcommand{\softmin}{\operatorname{softmin}}
\newcommand{\softmax}{\operatorname{softmax}}
\newcommand{\CVaR}{\mathrm{CVaR}}
\newcommand{\LogSumExp}{\operatorname{LogSumExp}}

\title{Differentiable Inverse Reachability for RM75:\\
Complete Mathematical Design of the IRDPLAYGROUND Operator}
\author{IRDPLAYGROUND technical note}
\date{}

\begin{document}
\maketitle

\begin{abstract}
This note gives a self-contained mathematical account of the inverse
reachability (IRD) operator implemented in \texttt{ird\_playground}.
The operator answers a narrower question than classical discrete IRD tables:
given a medically valid tool-centre-point (TCP) pose and a rail-induced robot
base (J1-axis) frame, what smooth direction increases collision-checked arm
reachability, optionally under a local position/orientation region?
We relate the design to four foundational papers
(Vahrenkamp et~al.\ 2013; Rudorfer 2024/RM4D; Murooka et~al.\ 2025; Chiu 1988),
derive every aggregation formula used in code, and explain why automatic
differentiation yields usable IRD gradients, why conditional queries over
local pose cones are well-defined, and why the same field supports
trajectory-parameter updates without retraining.
\end{abstract}

\tableofcontents
\newpage

% -----------------------------------------------------------------------------
\section{Problem statement and classical IRD}
\label{sec:problem}

\subsection{Reachability and inverse reachability}
\label{sec:classical}

Let the robot configuration be \(q\in\mathcal{Q}\subset\R^{n}\) and let
forward kinematics map configurations to TCP poses
\begin{equation}
\label{eq:fk}
  \mathrm{FK}:\mathcal{Q}\to\SE(3),\qquad
  q\mapsto \Tcp(q).
\end{equation}
Here \(\SE(3)\) is the special Euclidean group of rigid poses
(rotation \(R\in\SO(3)\) and translation \(p\in\R^{3}\)).

\begin{definition}[Binary reachability]
A pose \(r\in\SE(3)\) is \emph{reachable} if there exists a collision-free
configuration \(q\) with \(\mathrm{FK}(q)=r\) (within IK tolerance) and joint
limits satisfied. Write
\begin{equation}
\label{eq:binary-reach}
  \mathcal{R}
  \;=\;
  \bigl\{\, r\in\SE(3)
  \;\big|\;
  \exists\, q\in\mathcal{Q}_{\mathrm{free}}:
  \mathrm{FK}(q)=r
  \,\bigr\}.
\end{equation}
\end{definition}

A classical \emph{reachability distribution} (RD) of Vahrenkamp et~al.\
(2013) stores, in a voxelisation of workspace relative to a fixed robot base,
a quality score such as the measure of joint configurations that map into
that voxel. An \emph{inverse reachability distribution} (IRD) inverts the same
data: given a fixed TCP (e.g.\ a grasp), it yields a distribution over
candidate base poses for which that TCP would be reachable.

Formally, if \(t_{i}=\Tcp_{i}^{-1}\) denotes a base-to-TCP sample stored in the
RD, then for a world TCP \(p\in\SE(3)\) the corresponding candidate base pose
is
\begin{equation}
\label{eq:ird-inversion}
  b_{i}(p) \;=\; p\, t_{i}^{-1}
  \;\in\;\SE(3)
\end{equation}
(with optional projection onto \(\SE(2)\) for planar base placement).
The IRD score of \(b_{i}(p)\) is inherited from the RD quality of \(t_{i}\).
This is an offline discrete table lookup.

\subsection{What IRDPLAYGROUND computes instead}
\label{sec:our-question}

Our first-stage operator does \emph{not} replace multi-seed IK or final
collision validation. It answers:

\begin{quote}
For a medically valid TCP pose \(\Tcp\) and a rail-induced J1-axis frame
\(\Taxis\), what smooth direction increases collision-checked arm
reachability clearance, including under a local position/orientation
uncertainty or task region?
\end{quote}

Denote the decision variables of interest by \(x\) (rail coordinate, TCP yaw
on a vessel surface, trajectory control points, \ldots). We construct a scalar
\begin{equation}
\label{eq:operator-goal}
  C(x)
  \;=\;
  \mathcal{A}\bigl[\, f\bigl(c(\Tcp(x),\Taxis(x))\bigr)\,\bigr]
\end{equation}
where:
\begin{itemize}[leftmargin=1.4em]
  \item \(c(\cdot)\) is a smooth RM4D-style chart (Section~\ref{sec:chart});
  \item \(f\) is a learned signed clearance field (Section~\ref{sec:field});
  \item \(\mathcal{A}\) is a differentiable set aggregator
        (Region~A, TaskCone, SetQuery, or Trajectory;
        Sections~\ref{sec:region}--\ref{sec:traj}).
\end{itemize}
Gradient-based updates use \(\nabla_{x} C\) from automatic differentiation
through the entire Torch graph.

% -----------------------------------------------------------------------------
\section{Paper map: what each reference contributes}
\label{sec:papers}

\begin{table}[h]
\centering
\begin{tabular}{@{}llp{6.6cm}@{}}
\toprule
Paper & Core idea absorbed & Where it appears \\
\midrule
Vahrenkamp et~al.\ 2013
  & IRD = fix TCP, score base poses
  & \(\Tcp\) fixed (or slowly varying); query via \(\Taxis(y_{\mathrm{rail}})\) \\
Rudorfer RM4D 2024
  & Quotient common yaw/roll symmetries $\to\sim$4D map shared by RM \& IRD
  & Flange chart \(c\); only base yaw quotiented for probe45 \\
Murooka et~al.\ 2025
  & Differentiable reachability map \(f_{R}(r)\) for optimisation;
    classifiers can have flat exterior gradients
  & Single-output signed clearance; boundary stencils; AD through queries \\
Chiu 1988
  & Postures should match task transmission requirements
  & Tip$\times$roll / uncertainty cones encode task-compatible free or
    uncertain sets inside \(C\) \\
\bottomrule
\end{tabular}
\end{table}

\paragraph{Murooka's differentiable reachability map.}
Murooka et~al.\ define a scalar \(f_{R}(r)\) with
\begin{equation}
\label{eq:murooka}
  f_{R}(r)\ge 0 \text{ if reachable},\qquad
  f_{R}(r)< 0 \text{ otherwise},
\end{equation}
trained as a smooth binary classifier (e.g.\ softplus logistic loss). When
\(f_{R}\) is continuous and differentiable w.r.t.\ task coordinates, the
gradient \(\partial f_{R}/\partial r\) can enter trajectory optimisation
without repeated IK. They also note that a pure decision function may become
nearly constant far outside the reachable set, yielding uninformative
gradients. Our design follows the same sign convention but supervises an
explicit signed clearance with local collision-checked stencils and boundary
normals, so that \(\nabla f\) remains informative near the frontier
(Section~\ref{sec:training}).

\paragraph{RM4D symmetry.}
Rudorfer observes that for many 6-/7-axis arms, reachability is approximately
invariant to base yaw about gravity and to TCP roll about the last joint,
reducing an \(\SE(3)\) query to an intrinsic 4-D quotient. We adopt the
dimensionality-reduction spirit but enforce only the symmetries that hold for
RM75 with the probe45 tool: base (J1) yaw is a true free symmetry; TCP/flange
roll is \emph{not}, because the TCP \(z\)-axis is offset by \(\sim 50^{\circ}\)
from the J7 axis.

\paragraph{Chiu's task compatibility.}
Chiu measures how well a posture's velocity/force ellipsoids align with a
task. We do not compute Yoshikawa indices online; instead we encode task
freedom (large tip$\times$roll cone) and medical uncertainty (tiny box$\times$cone)
as set queries over \(f\), so that \(C\) is a task-conditioned reachability
margin.

% -----------------------------------------------------------------------------
\section{Frames, rail base, and the flange chart}
\label{sec:chart}

\subsection{World poses and relative pose}
\label{sec:frames}

A homogeneous pose is written
\begin{equation}
\label{eq:pose-matrix}
  T
  \;=\;
  \begin{bmatrix}
    R & p \\
    0^{\mathsf T} & 1
  \end{bmatrix}
  \in\SE(3),
  \qquad
  R\in\SO(3),\; p\in\R^{3}.
\end{equation}
Given world TCP \(\Tcp^{W}\) and world J1-axis frame \(\Taxis^{W}\), the TCP
expressed in the axis frame is
\begin{equation}
\label{eq:relative-pose}
\begin{aligned}
  R_{\mathrm{axis\leftarrow tcp}}
  &=
  (R_{\mathrm{axis}}^{W})^{\mathsf T}\, R_{\mathrm{tcp}}^{W},
  \\
  p_{\mathrm{axis\leftarrow tcp}}
  &=
  (R_{\mathrm{axis}}^{W})^{\mathsf T}
  \bigl(p_{\mathrm{tcp}}^{W}-p_{\mathrm{axis}}^{W}\bigr).
\end{aligned}
\end{equation}
\emph{Interpretation.}
Equation~\eqref{eq:relative-pose} is the group product
\((\Taxis^{W})^{-1}\Tcp^{W}\). All reachability queries are posed in this
relative frame, which is exactly the IRD viewpoint: move the base/axis and
the relative TCP changes even if the world TCP is fixed.

\subsection{Rail-induced axis frame}
\label{sec:rail}

For a linear rail along local axis \(e_{a}\) (code default \(a=1\), i.e.\ \(y\)),
\begin{equation}
\label{eq:rail-base}
  \Taxis^{W}(y)
  \;=\;
  T_{\mathrm{world\leftarrow rail}}\;
  \underbrace{
    \begin{bmatrix}
      I_{3} & y\, e_{a} \\
      0^{\mathsf T} & 1
    \end{bmatrix}
  }_{\text{pure translation by rail coordinate }y}\;
  T_{\mathrm{rail\leftarrow axis}}(y{=}0).
\end{equation}
\emph{Interpretation.}
\(y\) is the raw URDF rail coordinate. Differentiating
\eqref{eq:rail-base} w.r.t.\ \(y\) moves only the axis origin along the rail;
rotation is constant. This is the primitive that turns ``IRD w.r.t.\ base
placement'' into ``IRD w.r.t.\ one scalar rail DoF''.

\subsection{TCP to flange}
\label{sec:tcp-flange}

Let \(T_{\mathrm{flange\leftarrow tcp}}\) be the fixed tool transform from
URDF. Flange (link\_7) pose in the axis frame is
\begin{equation}
\label{eq:tcp-to-flange}
  \bigl(p_{\mathrm{fl}},R_{\mathrm{fl}}\bigr)
  \;=\;
  \mathrm{compose}\bigl(
    (p_{\mathrm{tcp}},R_{\mathrm{tcp}}),\;
    T_{\mathrm{flange\leftarrow tcp}}
  \bigr).
\end{equation}
We score flange geometry rather than TCP geometry so that the chart axes
align with the J7 rotation axis (needed for the invariant / roll split below).

\subsection{Nine-dimensional yaw-invariant flange embedding}
\label{sec:9d}

Write flange position \(p=(p_{x},p_{y},p_{z})\) and flange axes
\(u_{x},u_{y},u_{z}\) (columns of \(R_{\mathrm{fl}}\)). Define the radial
smoothed distance
\begin{equation}
\label{eq:radial}
  r
  \;=\;
  \sqrt{p_{x}^{2}+p_{y}^{2}+\varepsilon},
  \qquad
  \varepsilon=(10^{-3}\,\mathrm{m})^{2}.
\end{equation}
The chart is the map
\begin{equation}
\label{eq:chart}
  c(p,R)
  \;=\;
  \begin{pmatrix}
    p_{z}\\
    r\\
    u_{x}\!\cdot\!\hat z\\
    p_{xy}\!\cdot\! u_{x,xy}\\
    p_{xy}\!\times\! u_{x,xy}\\
    u_{z}\!\cdot\!\hat z\\
    p_{xy}\!\cdot\! u_{z,xy}\\
    p_{xy}\!\times\! u_{z,xy}\\
    u_{y}\!\cdot\!\hat z
  \end{pmatrix}
  \in\R^{9},
\end{equation}
where \(p_{xy}=(p_{x},p_{y})\), \(u_{\cdot,xy}\) is the \(xy\)-projection of an
axis, \(\hat z=(0,0,1)\), the 2-D ``cross'' is the scalar
\(p_{x}u_{y}-p_{y}u_{x}\), and
\begin{equation}
\label{eq:uyz}
  u_{y}\!\cdot\!\hat z
  \;=\;
  (u_{z}\times u_{x})\!\cdot\!\hat z
  \;=\;
  u_{z,x}\,u_{x,y}-u_{z,y}\,u_{x,x}.
\end{equation}

\paragraph{Why each component.}
\begin{itemize}[leftmargin=1.4em]
  \item \(p_{z}\): height along the J1 axis (invariant to yaw about \(\hat z\)).
  \item \(r\): cylindrical radius; \(\varepsilon\) removes the kink of
        \(\sqrt{p_{x}^{2}+p_{y}^{2}}\) at the axis (curvature \(\sim 10^{3}/\mathrm{m}\)
        instead of \(\sim 10^{6}/\mathrm{m}\)).
  \item Components involving \(u_{z}\): J7-invariant block
        (indices \(\{0,1,5,6,7\}\) in code), because spinning about J7 rotates
        \(u_{x},u_{y}\) but not \(u_{z}\).
  \item Components involving \(u_{x}\) and \(u_{y}\!\cdot\!\hat z\):
        \emph{roll-carrying} block. They retain flange roll \(\gamma\), which
        is a genuine DoF for probe45 and must not be quotiented.
  \item The ninth entry \(u_{y}\!\cdot\!\hat z\) completes \(R^{\mathsf T}\hat z\).
        Without it, as \(r\to 0\) the mixed products vanish and a twofold
        ambiguity remains in reconstructing orientation about the axis.
\end{itemize}

\paragraph{Symmetry statement.}
If \(R_{z}(\alpha)\) is a rotation about the J1 axis applied simultaneously to
base and (equivalently) to the relative flange pose, then
\begin{equation}
\label{eq:yaw-invariance}
  c\bigl(R_{z}(\alpha)\,p,\; R_{z}(\alpha)\,R\bigr)
  \;=\;
  c(p,R).
\end{equation}
Thus \(c\) embeds the intrinsic quotient
\(\SE(3)/\mathrm{Yaw}_{J1}\) smoothly into \(\R^{9}\) (redundant coordinates
of an intrinsic 5-D flange chart after quotienting only yaw).

\begin{remark}[Legacy 5-D TCP chart]
An older embedding
\(c_{\mathrm{TCP}}=[p_{z},\,u_{z}\!\cdot\!\hat z,\,\|p_{xy}\|,\,
p_{xy}\!\cdot\!u_{xy},\,p_{xy}\times u_{xy}]\)
quotiented both yaw and TCP roll. For probe45 this is incorrect and is kept
only for audits.
\end{remark}

% -----------------------------------------------------------------------------
\section{Signed reachability field}
\label{sec:field}

\subsection{Definition}
\label{sec:field-def}

\begin{definition}[Signed clearance field]
The network \(f_{\theta}:\R^{9}\to\R\) approximates a signed reachability
clearance on the flange chart:
\begin{equation}
\label{eq:sign-convention}
  f_{\theta}(c)
  \;
  \begin{cases}
    >0 & \text{reachable (positive margin)},\\
    =0 & \text{approximate boundary},\\
    <0 & \text{unreachable}.
  \end{cases}
\end{equation}
\end{definition}
This matches Murooka's sign pattern~\eqref{eq:murooka} but the \emph{magnitude}
is trained to track a collision-checked local distance in a declared metric
(Section~\ref{sec:metric}), not merely a classification margin.

\subsection{Input normalisation and Fourier features}
\label{sec:fourier}

Training chart rows determine centre \(\mu\in\R^{9}\) and scale \(\sigma\in\R^{9}\)
(quantile mid-range). Normalised inputs are
\begin{equation}
\label{eq:normalize}
  \tilde c
  \;=\;
  \frac{c-\mu}{\sigma}
  \qquad
  \text{(elementwise)}.
\end{equation}
Fourier features with \(B\) bands (\(B{=}3\) by default) are
\begin{equation}
\label{eq:fourier}
  \Phi(\tilde c)
  \;=\;
  \Bigl[
    \tilde c,\;
    \sin(\pi 2^{0}\tilde c),\;
    \cos(\pi 2^{0}\tilde c),\;
    \ldots,\;
    \sin(\pi 2^{B-1}\tilde c),\;
    \cos(\pi 2^{B-1}\tilde c)
  \Bigr]
  \in\R^{9(1+2B)}.
\end{equation}
\emph{Interpretation.}
Metres and dimensionless orientation channels have incompatible dynamic
ranges; without~\eqref{eq:normalize}, \(\sin(\pi 2^{k}x)\) would scramble
them. Fourier lifting helps represent sharp reachable/unreachable transitions
without a spatial hash grid (the previous model memorised boundary samples).

\subsection{Network architecture}
\label{sec:mlp}

Let \(\mathrm{softplus}_{\beta}(z)=\beta^{-1}\log(1+e^{\beta z})\). The field is
a residual MLP with Softplus activations and a linear scalar head:
\begin{equation}
\label{eq:mlp}
\begin{aligned}
  h_{0}
  &=
  \mathrm{softplus}_{\beta}\bigl(W_{0}\Phi(\tilde c)+b_{0}\bigr),
  \\
  h_{\ell+1}
  &=
  \mathrm{softplus}_{\beta}\bigl(
    h_{\ell}+W_{\ell}^{(2)}\,
    \mathrm{softplus}_{\beta}(W_{\ell}^{(1)}h_{\ell}+b_{\ell}^{(1)})
    +b_{\ell}^{(2)}
  \bigr),
  \\
  f_{\theta}(c)
  &=
  w^{\mathsf T}h_{L}+b.
\end{aligned}
\end{equation}
Softplus (unlike ReLU) is \(C^{\infty}\) with nowhere-vanishing derivative for
finite inputs, which is desirable when \(\nabla_{c}f\) is used as a guidance
direction.

\subsection{Support guard}
\label{sec:guard}

Let \([\underline c,\overline c]\) be the axis-aligned box of training chart
rows. After normalisation, a smooth exterior distance is
\begin{equation}
\label{eq:support-dist}
  d_{\mathrm{supp}}(c)
  \;=\;
  \sqrt{
    \bigl\|
      \mathrm{ReLU}(\underline{\tilde c}-\tilde c)
      +
      \mathrm{ReLU}(\tilde c-\overline{\tilde c})
    \bigr\|_{2}^{2}
    +\epsilon^{2}
  }
  -\epsilon,
\end{equation}
and the guarded score is
\begin{equation}
\label{eq:guarded}
  \tilde f(c)
  \;=\;
  f_{\theta}(c)
  -
  \lambda_{g}\, d_{\mathrm{supp}}(c),
  \qquad
  \lambda_{g}=8\text{ (default)}.
\end{equation}
\emph{Interpretation.}
Outside the training support, \(\tilde f\) is pulled negative so optimisation
is discouraged from extrapolating into unlabelled chart regions.

\subsection{World scoring (the IRD point query)}
\label{sec:score-world}

\begin{equation}
\label{eq:score-world}
  F(\Tcp^{W},\Taxis^{W})
  \;=\;
  \tilde f\Bigl(
    c\bigl(
      \mathrm{flange}\bigl(
        (\Taxis^{W})^{-1}\Tcp^{W}
      \bigr)
    \bigr)
  \Bigr).
\end{equation}
\emph{Interpretation.}
Fix \(\Tcp^{W}\) and vary \(\Taxis^{W}(y)\) via~\eqref{eq:rail-base}: then
\(F\) is precisely a smooth IRD score of the rail/base placement. Fix the
axis and vary \(\Tcp^{W}\): then \(F\) is a smooth forward reachability score.
One field serves both views because of the relative chart~\eqref{eq:chart}.

% -----------------------------------------------------------------------------
\section{Declared SE(3) metric and training targets}
\label{sec:training}

\subsection{Metric}
\label{sec:metric}

Boundary labels use the weighted distance that equates \(1\,\mathrm{cm}\) of
translation with \(1^{\circ}\) of rotation:
\begin{equation}
\label{eq:lambda-metric}
  \lambda
  \;=\;
  \frac{0.01\,\mathrm{m}}{1^{\circ}\text{ in radians}}
  \;=\;
  \frac{0.01}{\pi/180}
  \;\approx\;
  0.5730\,\mathrm{m/rad}.
\end{equation}
For a translational increment \(\Delta p\in\R^{3}\) (metres) and rotational
increment \(\Delta\omega\in\R^{3}\) (radians, rotation-vector / so(3) coordinates),
\begin{equation}
\label{eq:se3-distance}
  d(\Delta p,\Delta\omega)
  \;=\;
  \sqrt{
    \|\Delta p\|_{2}^{2}
    +
    \lambda^{2}\,\|\Delta\omega\|_{2}^{2}
  }.
\end{equation}
\emph{Interpretation.}
This is \emph{not} the bi-invariant \(\SE(3)\) geodesic; it is a declared
chart-friendly metric used to place stencil samples and to interpret
clearance in ``metres''.

\subsection{Stencil supervision}
\label{sec:stencil}

Around collision-checked boundary centres one evaluates exact reachable /
unreachable labels at offsets
\begin{equation}
\label{eq:stencil-offsets}
  \pm 1,3,6,10\,\mathrm{mm}
  \quad\text{and}\quad
  \pm 1,3,5^{\circ},
\end{equation}
plus explicit zero-valued bisected boundary centres. These provide targets
\(s^{\star}(c)\) for the signed value of \(f_{\theta}(c)\).

\subsection{Losses}
\label{sec:losses}

Let \(\{(c_{i},y_{i},s_{i}^{\star},n_{i})\}_{i}\) be a minibatch with binary
label \(y_{i}\in\{0,1\}\), SDF target \(s_{i}^{\star}\), and (when available)
unit boundary normal \(n_{i}\) in normalised chart coordinates.

\paragraph{Classification (sign).}
\begin{equation}
\label{eq:loss-cls}
  \mathcal{L}_{\mathrm{cls}}
  \;=\;
  \frac{
    \sum_{i} w_{i}^{\mathrm{cls}}\,
    \ell_{\mathrm{BCE}}\bigl(\alpha f_{\theta}(c_{i}),\, y_{i}\bigr)
  }{
    \sum_{i} w_{i}^{\mathrm{cls}}
  },
\end{equation}
where \(\alpha\) is a logit scale (default \(3\)) and
\(\ell_{\mathrm{BCE}}\) is binary cross-entropy with logits.
\emph{Interpretation.}
Enforces Murooka-style sign correctness~\eqref{eq:sign-convention}.

\paragraph{Signed value (magnitude).}
On rows with SDF weight \(w_{i}^{\mathrm{sdf}}>0\),
\begin{equation}
\label{eq:loss-sdf}
  \mathcal{L}_{\mathrm{sdf}}
  \;=\;
  \mathrm{SmoothL1}_{\beta=0.1}
  \bigl(f_{\theta}(c_{i}),\, s_{i}^{\star}\bigr).
\end{equation}
\emph{Interpretation.}
Makes \(|f|\) meaningful as clearance, not an arbitrary classifier margin.

\paragraph{Boundary normal alignment.}
With \(\tilde c\) requiring grad,
\begin{equation}
\label{eq:loss-normal}
  \mathcal{L}_{\mathrm{n}}
  \;=\;
  \mathbb{E}_{i}\Bigl[
    1
    -
    \frac{
      \nabla_{\tilde c}f_{\theta}(c_{i})
      \;\cdot\;
      n_{i}
    }{
      \|\nabla_{\tilde c}f_{\theta}(c_{i})\|\,\|n_{i}\|
    }
  \Bigr].
\end{equation}
\emph{Interpretation.}
Cosine loss aligns the field gradient with the collision-checked crossing
direction. This is why \(\nabla F\) points toward increasing reachability
near the frontier rather than along arbitrary level sets.

\paragraph{Optional Eikonal (disabled in production).}
\begin{equation}
\label{eq:loss-eikonal}
  \mathcal{L}_{\mathrm{Eik}}
  \;=\;
  \mathrm{SmoothL1}\bigl(
    \|\nabla_{\tilde c}f\|,\;
    v^{\star}
  \bigr),
\end{equation}
with \(v^{\star}=1\) or an empirical stencil slope. Workspace-normalised
mixed units make a universal unit-speed prior poorly conditioned; production
relies on~\eqref{eq:loss-sdf}--\eqref{eq:loss-normal} instead.

\paragraph{Total.}
\begin{equation}
\label{eq:loss-total}
  \mathcal{L}
  \;=\;
  \lambda_{\mathrm{cls}}\mathcal{L}_{\mathrm{cls}}
  +
  \lambda_{\mathrm{sdf}}\mathcal{L}_{\mathrm{sdf}}
  +
  \lambda_{\mathrm{n}}\mathcal{L}_{\mathrm{n}}
  +
  \lambda_{\mathrm{Eik}}\mathcal{L}_{\mathrm{Eik}}
  +
  \cdots
\end{equation}

\begin{remark}[Why not a classifier alone?]
A classifier can achieve high accuracy while being nearly constant for
\(c\) deep in the unreachable region, so \(\nabla_{x}F\approx 0\) exactly
when an optimiser most needs a recovery direction (Murooka's caveat). Signed
stencil values plus normal alignment keep a usable slope across the boundary
band.
\end{remark}

% -----------------------------------------------------------------------------
\section{Differentiable set queries}
\label{sec:region}

All set queries share the pattern: sample poses about a nominal TCP, score
each with~\eqref{eq:score-world}, reduce with a smooth aggregator.

\subsection{Local pose perturbation}
\label{sec:perturb}

In the medical TCP frame with columns \([b,t,n]\) (binormal, tangent, normal),
a position offset \(\delta p_{\mathrm{loc}}\in\R^{3}\) and rotation vector
\(\delta\omega_{\mathrm{loc}}\in\R^{3}\) produce
\begin{equation}
\label{eq:perturb}
\begin{aligned}
  p'
  &=
  p + R\,\delta p_{\mathrm{loc}},
  \\
  R'
  &=
  R\;\operatorname{Exp}(\delta\omega_{\mathrm{loc}}),
\end{aligned}
\end{equation}
where \(\operatorname{Exp}:\R^{3}\to\SO(3)\) is the Rodrigues map
\begin{equation}
\label{eq:rodrigues}
  \operatorname{Exp}(\omega)
  \;=\;
  I
  +
  \frac{\sin\theta}{\theta}[\hat\omega]
  +
  \frac{1-\cos\theta}{\theta^{2}}[\hat\omega]^{2},
  \qquad
  \theta=\|\omega\|,
\end{equation}
and \([\hat\omega]\) is the skew-symmetric matrix of the unit axis
(with the usual \(\theta\to 0\) limit \(I+[\omega]\)).

\subsection{Normalised softmin and softmax}
\label{sec:soft}

For a vector \(v\in\R^{K}\) and temperature \(\tau>0\),
\begin{equation}
\label{eq:logsumexp}
  \LogSumExp(v)
  \;=\;
  \log\sum_{k=1}^{K} e^{v_{k}}.
\end{equation}
The \emph{normalised softmin} used for robust aggregation is
\begin{equation}
\label{eq:softmin}
  \softmin_{\tau}(v)
  \;=\;
  -\tau\Bigl(
    \LogSumExp(-v/\tau)
    -
    \log K
  \Bigr).
\end{equation}
\emph{Why the \(\log K\) term?}
If \(v_{k}\equiv a\) for all \(k\), then
\(\LogSumExp(-v/\tau)=\log K - a/\tau\), hence
\(\softmin_{\tau}(v)=a\). A constant field is unchanged; without
\(\log K\), softmin would shift by \(-\tau\log K\).

The \emph{normalised softmax-like maximum} is
\begin{equation}
\label{eq:softmax}
  \softmax_{\tau}(v)
  \;=\;
  \tau\Bigl(
    \LogSumExp(v/\tau)
    -
    \log K
  \Bigr).
\end{equation}
Same constant-preserving property for maxima.

\paragraph{Implicit weights (gradient structure).}
Differentiating~\eqref{eq:softmin} gives
\begin{equation}
\label{eq:softmin-grad}
  \frac{\partial\softmin_{\tau}(v)}{\partial v_{k}}
  \;=\;
  \frac{e^{-v_{k}/\tau}}{\sum_{j}e^{-v_{j}/\tau}}
  \;=\;
  \pi_{k}^{\mathrm{min}}(v),
\end{equation}
i.e.\ a Gibbs distribution concentrated on the \emph{smallest} entries.
For~\eqref{eq:softmax},
\begin{equation}
\label{eq:softmax-grad}
  \frac{\partial\softmax_{\tau}(v)}{\partial v_{k}}
  \;=\;
  \frac{e^{v_{k}/\tau}}{\sum_{j}e^{v_{j}/\tau}}
  \;=\;
  \pi_{k}^{\mathrm{max}}(v),
\end{equation}
concentrated on the \emph{largest} entries. Thus robust queries backpropagate
into worst-case samples; task-free queries backpropagate into best-case
samples.

\subsection{Region A: uncertainty box \(\times\) small cone (softmin)}
\label{sec:region-a}

\begin{definition}[Region A]
Default extents: binormal \(4\,\mathrm{mm}\), tangent \(3\,\mathrm{mm}\),
normal \(2\,\mathrm{mm}\); orientation cone half-angle \(\beta=3^{\circ}\);
\(K=64\) fixed scrambled Sobol samples including the nominal pose.
\end{definition}

Position half-samples use a joint 5-D Sobol sequence \(u\in(0,1)^{5}\):
\begin{equation}
\label{eq:region-pos}
  \delta p
  \;=\;
  (2u_{1:3}-1)\odot
  (b,t,n)_{\mathrm{extent}},
\end{equation}
and orientation samples are uniform in solid angle inside the cone,
\begin{equation}
\label{eq:region-ori}
\begin{aligned}
  \cos\rho
  &=
  1-u_{4}\bigl(1-\cos\beta\bigr),
  \\
  \rho
  &=
  \arccos(\cos\rho),
  \\
  \phi
  &=
  2\pi u_{5},
  \\
  \delta\omega
  &=
  (\rho\cos\phi,\;\rho\sin\phi,\;0),
\end{aligned}
\end{equation}
then antisymmetrised \(\{\,0,\,\pm\text{half samples}\,\}\) to length \(K\).
The Region~A clearance is
\begin{equation}
\label{eq:region-a}
  C_{\mathrm{A}}(\Tcp,\Taxis)
  \;=\;
  \softmin_{\tau}
  \Bigl\{
    F\bigl(\Tcp(\delta_{k}),\,\Taxis\bigr)
  \Bigr\}_{k=1}^{K},
  \qquad
  \tau=0.15.
\end{equation}
\emph{Interpretation.}
\(C_{\mathrm{A}}\) is a smooth lower envelope over a medically motivated
uncertainty set \(U\). Increasing \(C_{\mathrm{A}}\) improves the worst
likely pose under small registration / contact error. Coverage diagnostics
use \(\frac{1}{K}\sum_{k}\sigma(F_{k}/\tau)\).

\subsection{Task cone: tip \(\times\) roll free set (softmax)}
\label{sec:task-cone}

Ultrasound contact often permits a large tip tilt and probe roll as
\emph{controller-chosen} free variables, not as uncertainty.

\begin{definition}[TaskCone]
Tip half-angle \(\beta_{\mathrm{tip}}\) (e.g.\ \(20^{\circ}\)--\(45^{\circ}\)),
roll half-range \(\gamma_{\max}\) (e.g.\ \(20^{\circ}\)--\(30^{\circ}\)),
\(K\) Sobol samples in \((\rho,\phi,\gamma)\).
\end{definition}
\begin{equation}
\label{eq:task-cone-sample}
\begin{aligned}
  \cos\rho
  &=
  1-u_{1}\bigl(1-\cos\beta_{\mathrm{tip}}\bigr),
  \\
  \phi
  &=
  2\pi u_{2},
  \\
  \gamma
  &=
  (2u_{3}-1)\gamma_{\max},
  \\
  \delta\omega
  &=
  (\rho\cos\phi,\;\rho\sin\phi,\;\gamma).
\end{aligned}
\end{equation}
Position is held fixed (contact point prescribed). The task clearance is
\begin{equation}
\label{eq:task-cone}
  C_{\mathrm{task}}(\Tcp,\Taxis)
  \;=\;
  \softmax_{\tau}
  \Bigl\{
    F\bigl(\Tcp(\eta_{k}),\,\Taxis\bigr)
  \Bigr\}_{k=1}^{K},
  \qquad
  \tau=0.20.
\end{equation}
\emph{Interpretation.}
Opposite of Region~A: \(C_{\mathrm{task}}\) asks whether \emph{some}
tip$\times$roll choice makes the arm comfortable. Gradients push the base/rail
toward placements that admit a high-clearance free pose inside the cone.
This is the Chiu-style ``match the posture set to the task'' idea, expressed
as a reachability margin rather than an ellipsoid index.

\subsection{Nested free \(\times\) uncertain query (SetQuery)}
\label{sec:set-query}

The general Phase-4 semantic is
\begin{equation}
\label{eq:nested-semantic}
  C_{\mathrm{set}}(T)
  \;=\;
  \max_{\eta\in D_{\mathrm{free}}}
  \;
  \underbrace{
    \CVaR_{\alpha}\text{ or }\softmin
  }_{\delta\in U}
  \;
  F\bigl(c(T(\eta,\delta))\bigr).
\end{equation}
Batch scores have shape \([N,K_{\mathrm{free}},K_{\mathrm{unc}}]\).
Reduce the uncertain axis first, then take a smooth max over free variables.

\paragraph{Rockafellar--Uryasev lower-tail CVaR.}
For clearance values (higher better), lower-tail CVaR at level \(\alpha\) is
\begin{equation}
\label{eq:cvar}
  \CVaR_{\alpha}(X)
  \;=\;
  \min_{t\in\R}
  \Biggl\{
    t
    +
    \frac{1}{1-\alpha}\,
    \mathbb{E}\bigl[(t-X)_{+}\bigr]
  \Biggr\}.
\end{equation}
\emph{Interpretation.}
\(t\) plays the role of a VaR threshold; the expectation penalises shortfalls
below \(t\). Empirically the optimum lies in the sample set, so code evaluates
the objective at every sample value and takes a hard or soft min over those
candidate \(t\)'s. This is more statistically stable than a pure softmin when
\(K_{\mathrm{unc}}\) is moderate.

\paragraph{Implemented reduction.}
\begin{equation}
\label{eq:set-query}
\begin{aligned}
  g(\eta)
  &=
  \CVaR_{\alpha}\bigl\{F(T(\eta,\delta_{j}))\bigr\}_{j}
  \quad\text{or}\quad
  \softmin_{\tau}\{\ldots\},
  \\
  C_{\mathrm{set}}
  &=
  \softmax_{\tau}\bigl\{g(\eta_{i})\bigr\}_{i}.
\end{aligned}
\end{equation}

% -----------------------------------------------------------------------------
\section{Why IRD gradients exist and are useful}
\label{sec:gradients}

\subsection{Chain rule through the IRD point query}
\label{sec:chain}

Consider rail coordinate \(y\) with fixed world TCP. Then
\begin{equation}
\label{eq:chain-rail}
  \frac{\mathrm{d}F}{\mathrm{d}y}
  \;=\;
  \frac{\partial\tilde f}{\partial c}\;
  \frac{\partial c}{\partial (p,R)}\;
  \frac{\partial (p,R)}{\partial \Taxis}\;
  \frac{\mathrm{d}\Taxis}{\mathrm{d}y}.
\end{equation}
Every factor is implemented with Torch ops (\(\operatorname{Exp}\), matrix
products, \(\sqrt{\,\cdot+\varepsilon}\), Softplus MLP), so
\(\mathrm{d}F/\mathrm{d}y\) is exact up to floating-point error. Empirical
AD-vs-FD median relative error on rail and TCP-\(x\) is on the order of
\(10^{-4}\).

\paragraph{Geometric reading.}
\(\partial c/\partial(p,R)\) converts a spatial twist of the relative flange
pose into chart coordinates. \(\partial\tilde f/\partial c\) is trained to
align with the outward reachable normal near the boundary
(Eq.~\eqref{eq:loss-normal}). Hence \(\mathrm{d}F/\mathrm{d}y>0\) means
``moving the carriage in the \(+y\) direction increases predicted clearance''.

\subsection{Gradients of conditional cone queries}
\label{sec:cond-grad}

For Region~A,
\begin{equation}
\label{eq:grad-region}
  \nabla_{x} C_{\mathrm{A}}
  \;=\;
  \sum_{k=1}^{K}
  \pi_{k}^{\mathrm{min}}\,
  \nabla_{x} F\bigl(\Tcp(\delta_{k};x),\,\Taxis(x)\bigr).
\end{equation}
For TaskCone,
\begin{equation}
\label{eq:grad-task}
  \nabla_{x} C_{\mathrm{task}}
  \;=\;
  \sum_{k=1}^{K}
  \pi_{k}^{\mathrm{max}}\,
  \nabla_{x} F\bigl(\Tcp(\eta_{k};x),\,\Taxis(x)\bigr).
\end{equation}
\emph{Interpretation of ``conditional IRD gradient''.}
The condition is the set \(U\) or \(D_{\mathrm{free}}\). The scalar \(C\) is
the set-aggregated clearance. Its gradient w.r.t.\ base/rail/path parameters
is a \emph{weighted average of point-IRD gradients} at the sampled local
poses, with weights~\eqref{eq:softmin-grad} or~\eqref{eq:softmax-grad}.
No discrete voxel argmax is required.

\subsection{Vessel-surface example \((\theta,y)\)}
\label{sec:vessel-grad}

On an elliptical vessel, TCP position is tied to surface angle \(\theta\) and
path abscissa \(s\), while rail \(y\) sets the base. Demo code forms
\begin{equation}
\label{eq:demo-query}
  C(s;\theta,y)
  \;=\;
  C_{\mathrm{task}}\bigl(
    \Tcp(\theta,s),\,
    \Taxis(y)
  \bigr)
\end{equation}
and extracts
\begin{equation}
\label{eq:demo-grads}
  \frac{\partial C}{\partial\theta},
  \qquad
  \frac{\partial C}{\partial y}
\end{equation}
by \(\mathtt{torch.autograd.grad}\). Quiver plots of these components are
literal visualisations of the conditional IRD vector field on the
\((\theta,y)\) planning plane.

% -----------------------------------------------------------------------------
\section{Trajectory operator and online updates}
\label{sec:traj}

\subsection{Aggregation order}
\label{sec:agg-order}

Given waypoints \(\{\Tcp^{(\ell)}\}_{\ell=1}^{W}\) and a (shared or
per-waypoint) axis frame, the trajectory operator applies, in order:
\begin{enumerate}[leftmargin=1.6em]
  \item \textbf{Uncertainty / Region~A scenarios} for each approximate-angle
        candidate;
  \item \textbf{Angle aggregation} (\(\softmax\) select, \(\softmin\) robust,
        mean, or CVaR);
  \item \textbf{Path residuals}: interior SE(3) samples on each segment;
  \item \textbf{Trajectory reduction}: softmin / CVaR / hard min over
        waypoint and path margins.
\end{enumerate}

\subsection{SE(3) path interpolation}
\label{sec:se3-interp}

On a segment from \(T_{0}\) to \(T_{1}\), for \(\alpha\in(0,1)\),
\begin{equation}
\label{eq:se3-interp}
\begin{aligned}
  p(\alpha)
  &=
  (1-\alpha)p_{0}+\alpha p_{1},
  \\
  R(\alpha)
  &=
  R_{0}\,
  \operatorname{Exp}\bigl(
    \alpha\,\operatorname{Log}(R_{0}^{\mathsf T}R_{1})
  \bigr).
\end{aligned}
\end{equation}
\emph{Interpretation.}
Translations move linearly; orientations follow the geodesic in \(\SO(3)\).
Controllers that require the whole interpolated path to be IK-feasible need
these interior samples: endpoint reachability alone is insufficient.

\subsection{Per-waypoint and path margins}
\label{sec:path-margin}

Let \(A\) index angle candidates and \(S\) index Region~A scenarios. After
scoring,
\begin{equation}
\label{eq:waypoint-clearance}
\begin{aligned}
  g_{\ell,a}
  &=
  \mathcal{A}_{S}\bigl[F(T^{(\ell)}_{a,s})\bigr],
  \\
  C_{\ell}
  &=
  \mathcal{A}_{A}\bigl[g_{\ell,a}\bigr],
\end{aligned}
\end{equation}
and for interior path samples \(C^{\mathrm{path}}_{m}\) similarly. The scalar
trajectory clearance is
\begin{equation}
\label{eq:traj-clearance}
  C_{\mathrm{traj}}
  \;=\;
  \softmin_{\tau}
  \Bigl[
    \{C_{\ell}\}_{\ell=1}^{W}
    \;\cup\;
    \{C^{\mathrm{path}}_{m}\}_{m}
  \Bigr]
\end{equation}
(or CVaR / hard min). Optional SRS arm angle \(\psi\) enters as an outer
\(\softmax\) over \(\psi\)-conditioned fields when exposed.

\subsection{Why the landscape updates with the trajectory}
\label{sec:why-update}

The network \(f_{\theta}\) is trained once in the flange chart and then frozen.
A trajectory is only a curve \(x\mapsto(\Tcp(x),\Taxis(x))\). Because:
\begin{enumerate}[leftmargin=1.6em]
  \item \(c\) depends on the \emph{relative} pose~\eqref{eq:relative-pose};
  \item cone samples are expressed in the \emph{current} TCP body frame
        \eqref{eq:perturb};
  \item aggregators~\eqref{eq:softmin}--\eqref{eq:softmax} are memoryless
        functions of the current sample scores;
\end{enumerate}
re-evaluating \(C_{\mathrm{task}}(s)\) or \(C_{\mathrm{traj}}(x)\) at a new
path automatically yields a new IRD landscape and new gradients
\(\nabla_{x}C\) \emph{without} modifying \(\theta\).

\subsection{Typical optimisation objective}
\label{sec:opt}

Ellipse-vessel demos minimise a hinge toward a conformal clearance target
\(m^{\star}\) plus regularisers:
\begin{equation}
\label{eq:opt-loss}
\begin{aligned}
  \ell_{\mathrm{IRD}}
  &=
  \frac{1}{W}\sum_{\ell}
  \mathrm{softplus}\!\Bigl(\frac{m^{\star}-C_{\ell}}{\sigma}\Bigr)
  +
  \max_{\ell}\mathrm{softplus}\!\Bigl(\frac{m^{\star}-C_{\ell}}{\sigma}\Bigr)
  -
  \kappa\,\overline{C},
  \\
  \ell
  &=
  w_{\mathrm{IRD}}\,\ell_{\mathrm{IRD}}
  +
  w_{\mathrm{near}}\,\ell_{\mathrm{nearest}}
  +
  w_{\mathrm{cont}}\,\ell_{\mathrm{continuity}}
  +
  \cdots
\end{aligned}
\end{equation}
\emph{Interpretation.}
\(\mathrm{softplus}((m^{\star}-C)/\sigma)\) is a smooth one-sided penalty:
inactive when \(C\ge m^{\star}\), rising when clearance falls below target.
The extra \(-\kappa\overline{C}\) term continues to reward climbing toward
IRD peaks even after the hinge is satisfied. Continuity penalties keep
\(\theta(s)\) and \(y(s)\) smooth. Gradient descent on~\eqref{eq:opt-loss}
is exactly ``update the trajectory using conditional IRD gradients''.

% -----------------------------------------------------------------------------
\section{End-to-end commutative diagram}
\label{sec:diagram}

\[
\begin{array}{c}
x
\;\xrightarrow{\;\text{kinematics / rail}\;}
(\Tcp(x),\Taxis(x))
\;\xrightarrow{\;\text{local samples}\;}
\{T_{k}(x)\}
\\[0.6em]
\xrightarrow{\;\eqref{eq:relative-pose}+\eqref{eq:chart}\;}
\{c_{k}\}
\;\xrightarrow{\;\tilde f\;}
\{F_{k}\}
\;\xrightarrow{\;\mathcal{A}\;}
C(x)
\;\xrightarrow{\;\nabla\;}
\nabla_{x}C
\end{array}
\]

\begin{proposition}[Single autograd tape]
If every arrow above is implemented with elementary differentiable maps, then
\(\nabla_{x}C\) exists almost everywhere and coincides with the analytic
chain rule. Discrete IRD voxel lookups break the tape at the table index;
the neural field does not.
\end{proposition}

% -----------------------------------------------------------------------------
\section{Scope, guarantees, and non-goals}
\label{sec:scope}

\begin{itemize}[leftmargin=1.4em]
  \item \textbf{Guidance, not certificate.}
        Held-out boundary errors are sub-millimetre when bracketed, but the
        field is an optimisation-direction operator. Hard feasibility still
        requires final multi-seed IK and collision checking.
  \item \textbf{Symmetry approximation.}
        Only J1 yaw is exact; residual probe-roll collision differences are
        rare and rejected downstream.
  \item \textbf{No branch / continuous \(\psi\) in the point field.}
        Offline labels are point reachability. Path residuals and continuous
        arm-angle selection live in the trajectory operator.
  \item \textbf{Not a full patient-trajectory planner by itself.}
        Region~A / TaskCone / TrajectoryTask are query layers over the point
        field; application demos add geometric path parameterisations and
        loss terms.
\end{itemize}

% -----------------------------------------------------------------------------
\section{Symbol table}
\label{sec:symbols}

\begin{center}
\begin{tabular}{@{}ll@{}}
\toprule
Symbol & Meaning \\
\midrule
\(\Tcp,\Taxis\) & World TCP pose; world J1-axis pose \\
\(c\in\R^{9}\) & Yaw-invariant flange chart \\
\(f_{\theta},\tilde f\) & Signed field; support-guarded field \\
\(F\) & World point score~\eqref{eq:score-world} \\
\(C_{\mathrm{A}},C_{\mathrm{task}},C_{\mathrm{set}},C_{\mathrm{traj}}\) &
  Aggregated clearances \\
\(y\) & Rail coordinate \\
\(\delta,\eta\) & Uncertainty sample; free (task) sample \\
\(\tau\) & Softmin/softmax temperature \\
\(\lambda\) & Translation--rotation metric scale~\eqref{eq:lambda-metric} \\
\(\pi^{\mathrm{min}},\pi^{\mathrm{max}}\) & Softmin / softmax weights \\
\bottomrule
\end{tabular}
\end{center}

% -----------------------------------------------------------------------------
\section{References}
\label{sec:refs}

\begin{enumerate}[leftmargin=1.6em]
  \item N.~Vahrenkamp, T.~Asfour, and R.~Dillmann,
        ``Robot placement based on reachability inversion,'' ICRA 2013.
  \item M.~Rudorfer,
        ``RM4D: A Combined Reachability and Inverse Reachability Map for
        Common 6-/7-axis Robot Arms by Dimensionality Reduction to 4D,''
        ICRA 2025 / arXiv:2410.06968.
  \item M.~Murooka et~al.,
        ``Learning Differentiable Reachability Maps for Optimization-based
        Humanoid Motion Generation,'' Humanoids 2025 / arXiv:2508.11275.
  \item S.~L.~Chiu,
        ``Task Compatibility of Manipulator Postures,''
        International Journal of Robotics Research, 1988.
  \item Internal design note:
        \texttt{ird\_playground/docs/rm4d\_signed\_ird\_design.md}.
\end{enumerate}

\appendix
\section{Code $\leftrightarrow$ equation map}
\label{sec:code-map}

\begin{center}
\begin{tabular}{@{}ll@{}}
\toprule
Equation & Primary module \\
\midrule
\eqref{eq:rail-base}
  & \texttt{region/operator.py::base\_from\_rail\_torch} \\
\eqref{eq:chart}
  & \texttt{ird/canonical.py::canonical\_flange\_invariants\_torch} \\
\eqref{eq:score-world}
  & \texttt{neural/signed\_field.py::score\_world} \\
\eqref{eq:softmin}--\eqref{eq:softmax}
  & \texttt{region/operator.py}, \texttt{region/set\_query.py} \\
\eqref{eq:region-a}
  & \texttt{region/operator.py::RegionA} \\
\eqref{eq:task-cone}
  & \texttt{region/task\_cone.py::TaskConeReachability} \\
\eqref{eq:nested-semantic}--\eqref{eq:set-query}
  & \texttt{region/set\_query.py::SetQueryOperator} \\
\eqref{eq:se3-interp}--\eqref{eq:traj-clearance}
  & \texttt{region/trajectory\_operator.py} \\
\eqref{eq:loss-cls}--\eqref{eq:loss-total}
  & \texttt{neural/train\_signed.py} \\
\eqref{eq:opt-loss}
  & \texttt{experiments/ellipse\_vessel\_ird\_demo.py} \\
\bottomrule
\end{tabular}
\end{center}

\end{document}
